\documentclass[11pt]{article}
\usepackage[margin=1in]{geometry}
\usepackage{amsmath,amssymb,amsthm}
\usepackage{mathtools}
\usepackage{hyperref}

\newtheorem{theorem}{Theorem}
\newtheorem{lemma}[theorem]{Lemma}
\theoremstyle{definition}
\newtheorem{definition}{Definition}
\newtheorem{example}{Example}

\title{Gradient and Jacobian}
\author{math-foundations / concept 14}
\date{}

\begin{document}
\maketitle

\subsection*{First, an example}
Consider the bowl-shaped scalar field $f(x, y) = x^2 + y^2$. Its partial
derivatives are $\partial f/\partial x = 2x$ and $\partial f/\partial y = 2y$,
so the gradient is the vector field
\[
\nabla f(x, y) = (2x,\, 2y).
\]
\emph{Read aloud:} ``nabla f at x, y equals the vector two-x, two-y.''

At the point $(3, 4)$ this gives $\nabla f(3, 4) = (6, 8)$, a vector of
length $\sqrt{6^2 + 8^2} = 10$ pointing radially outward from the origin.
This is the direction of steepest ascent: moving an infinitesimal step
along $(6, 8)$ increases $f$ faster than any other direction of equal
length, and the rate of that increase is $10$. Negating it gives the
direction of steepest descent\,---\,the move that gradient descent would
take from $(3, 4)$. The rest of this lesson develops the gradient and its
vector-valued cousin, the Jacobian, in full generality.

\section{Setup and motivation}

We work in finite-dimensional Euclidean space $\mathbb{R}^n$ with the standard
inner product $\langle u, v\rangle = u^\top v$ and induced norm
$\|v\| = \sqrt{\langle v, v\rangle}$. For scalar functions, the derivative
captures a single rate of change; for vector-valued functions, it captures
\emph{all} pairwise rates between input and output coordinates, organized
into a matrix.

\section{Definitions}

\begin{definition}[Partial derivative]
Let $U \subseteq \mathbb{R}^n$ be open and $f : U \to \mathbb{R}$. The
\emph{partial derivative} of $f$ with respect to $x_j$ at $x \in U$ is
\[
\frac{\partial f}{\partial x_j}(x)
   = \lim_{h \to 0} \frac{f(x + h e_j) - f(x)}{h},
\]
\emph{Read aloud:} ``partial f by partial x-j at x equals the limit, as h tends to zero, of f of x plus h times e-j minus f of x, all divided by h.''

where $e_j$ is the $j$th standard basis vector, when this limit exists.
\end{definition}

\begin{definition}[Gradient]
Let $f : U \subseteq \mathbb{R}^n \to \mathbb{R}$ be differentiable at $x$.
The \emph{gradient} of $f$ at $x$ is the column vector
\[
\nabla f(x) \;=\; \Bigl( \tfrac{\partial f}{\partial x_1}(x),\, \ldots,\,
\tfrac{\partial f}{\partial x_n}(x) \Bigr)^{\!\top} \in \mathbb{R}^n.
\]
\emph{Read aloud:} ``nabla f at x equals the column vector of partial f by partial x-1 through partial f by partial x-n, transposed, an element of R-to-the-n.''

Equivalently, $\nabla f(x)$ is the unique vector such that
$D_v f(x) = \langle \nabla f(x), v\rangle$ for every $v \in \mathbb{R}^n$,
where $D_v f(x) = \lim_{t\to 0}\bigl(f(x + tv) - f(x)\bigr)/t$ is the
directional derivative.
\end{definition}

\begin{definition}[Jacobian matrix]
Let $F : U \subseteq \mathbb{R}^n \to \mathbb{R}^m$ with components
$F = (F_1, \ldots, F_m)$, differentiable at $x$. The \emph{Jacobian
matrix} of $F$ at $x$ is the $m \times n$ matrix
\[
J_F(x) \;=\; \begin{pmatrix}
\dfrac{\partial F_1}{\partial x_1}(x) & \cdots & \dfrac{\partial F_1}{\partial x_n}(x) \\[1ex]
\vdots & \ddots & \vdots \\[1ex]
\dfrac{\partial F_m}{\partial x_1}(x) & \cdots & \dfrac{\partial F_m}{\partial x_n}(x)
\end{pmatrix},
\qquad [J_F(x)]_{ij} = \frac{\partial F_i}{\partial x_j}(x).
\]
\emph{Read aloud:} ``the i-j entry of the Jacobian J-F at x is partial F-i by partial x-j.''

$J_F(x)$ is the matrix of the total derivative
$dF_x : \mathbb{R}^n \to \mathbb{R}^m$ in standard bases.
\end{definition}

\begin{definition}[Jacobian determinant]
If $m = n$, the \emph{Jacobian determinant} of $F$ at $x$ is
$\det J_F(x) \in \mathbb{R}$. It is the signed factor by which $F$
locally scales $n$-dimensional volume at $x$: a small parallelepiped
of volume $V$ at $x$ has image of volume $|\det J_F(x)|\,V + o(V)$.
\end{definition}

\section{The gradient and steepest ascent}

\begin{lemma}[Directional-derivative formula]
\label{lem:dirderiv}
If $f$ is differentiable at $x$, then for every $v \in \mathbb{R}^n$,
\[
D_v f(x) = \langle \nabla f(x), v\rangle.
\]
\emph{Read aloud:} ``the directional derivative D-v of f at x equals the inner product of nabla f at x with v.''
\end{lemma}
\begin{proof}[Sketch]
Differentiability at $x$ means $f(x + h) = f(x) + \langle \nabla f(x), h\rangle + r(h)$
with $r(h)/\|h\| \to 0$. Setting $h = tv$ and dividing by $t$ gives the claim
in the limit $t \to 0$.
\end{proof}

\begin{theorem}[Direction of maximum directional derivative]
\label{thm:steepest}
Let $f : U \subseteq \mathbb{R}^n \to \mathbb{R}$ be differentiable at
$x \in U$, with $\nabla f(x) \neq 0$. Then among all unit vectors
$v \in \mathbb{R}^n$ (i.e.\ $\|v\| = 1$), the directional derivative
$D_v f(x)$ attains its maximum at
\[
v^\star = \frac{\nabla f(x)}{\|\nabla f(x)\|},
\]
\emph{Read aloud:} ``v-star equals nabla f at x divided by the norm of nabla f at x.''

and the maximum value is $\|\nabla f(x)\|$.
\end{theorem}

\begin{proof}
By Lemma~\ref{lem:dirderiv}, $D_v f(x) = \langle \nabla f(x), v\rangle$.
The Cauchy--Schwarz inequality states that for any $u, v \in \mathbb{R}^n$,
\[
|\langle u, v\rangle| \;\le\; \|u\|\,\|v\|,
\]
\emph{Read aloud:} ``the absolute value of the inner product of u and v is at most the norm of u times the norm of v.''

with equality iff $u$ and $v$ are linearly dependent. Applying this with
$u = \nabla f(x)$ and a unit vector $v$,
\[
D_v f(x) = \langle \nabla f(x), v\rangle \;\le\; \|\nabla f(x)\|\,\|v\|
        = \|\nabla f(x)\|.
\]
\emph{Read aloud:} ``D-v f at x equals the inner product of nabla f at x with v, which is at most the norm of nabla f at x times the norm of v, equal to the norm of nabla f at x.''

Equality holds iff $v$ is a positive scalar multiple of $\nabla f(x)$;
since $\|v\| = 1$ and $\nabla f(x) \neq 0$, the unique maximizer is
$v^\star = \nabla f(x)/\|\nabla f(x)\|$. Substituting back gives
$D_{v^\star} f(x) = \langle \nabla f(x), \nabla f(x)/\|\nabla f(x)\|\rangle
= \|\nabla f(x)\|$. \qedhere
\end{proof}

\begin{example}
For $f(x, y) = x^2 + 3y^2$, $\nabla f(x, y) = (2x, 6y)$. At $(1, 1)$,
the steepest-ascent direction is $(2, 6)/\sqrt{40} = (1, 3)/\sqrt{10}$
with rate $\sqrt{40}$, while $-\nabla f$ points toward the minimum at
the origin.
\end{example}

\section{The chain rule for vector-valued maps}

\begin{theorem}[Chain rule]
Let $G : \mathbb{R}^p \to \mathbb{R}^n$ be differentiable at $x$, and
$F : \mathbb{R}^n \to \mathbb{R}^m$ differentiable at $y = G(x)$.
Then $H = F \circ G$ is differentiable at $x$ and
\[
J_H(x) \;=\; J_F(G(x)) \cdot J_G(x),
\]
\emph{Read aloud:} ``the Jacobian of H at x equals the Jacobian of F at G of x, times the Jacobian of G at x.''

an $m \times p$ matrix obtained as the product of an $m \times n$ matrix
and an $n \times p$ matrix.
\end{theorem}

This is the algebraic backbone of reverse-mode automatic differentiation:
for a deep composition $F_L \circ F_{L-1} \circ \cdots \circ F_1$, the
overall Jacobian factors as a product of layer Jacobians, and the loss
gradient propagates by successive vector--Jacobian products.

\section{Worked example: $F(x, y) = (x^2 - y, e^{xy})$}

The partials are $\partial F_1/\partial x = 2x$, $\partial F_1/\partial y = -1$,
$\partial F_2/\partial x = y e^{xy}$, $\partial F_2/\partial y = x e^{xy}$.
Hence
\[
J_F(x, y) = \begin{pmatrix} 2x & -1 \\ y e^{xy} & x e^{xy} \end{pmatrix},
\qquad
J_F(1, 0) = \begin{pmatrix} 2 & -1 \\ 0 & 1 \end{pmatrix},
\quad \det J_F(1, 0) = 2.
\]
\emph{Read aloud:} ``J-F at x, y is the two-by-two matrix with rows two-x, minus one and y-e-to-the-x-y, x-e-to-the-x-y; at one-zero this becomes two, minus one, zero, one, with determinant two.''

By the inverse function theorem, $F$ is a local diffeomorphism near $(1, 0)$,
locally doubling area.

\end{document}
